% ==============================================================================
% T0 Theory / FFGFT: Document 276 (Chapter Content, EN)
% Title: Homology test of the HLV vortex hypothesis -- cell-shape-independent
%        b1 filtration, deformation null test, and N-scaling
% Author: Johann Pascher
% Date: June 2026
% References: Doc. 270, 271, 274, 275; Krueger (HLV) Zenodo 20668125, 20629842, 20615649
% Companion script: Dok276_Skripte/ffgft_hlv_homology_check_v3.py
% ==============================================================================

\newtheorem*{observation*}{Observation}

\section*{Abstract}
Marcel Krüger's HLV G-Lattice programme~\cite{krueger276main} formulates a
finite 6D$\to$3D cut-and-project construction and tests it through a ladder of
pre-registered falsification runs (A--K7 plus the D1--D3 cell-geometry
addendum). The D1--D3 audit \emph{withdrew} the original claim of exact
dodecahedral Voronoi cells and replaced it with a narrower profile: a
non-crystalline, mixed-face, $\varphi$-skewed cell ensemble. The revised
vortex axiom therefore no longer requires a Platonic cell, but closed cycles,
phase continuity, and holonomy-carrying loops. Exactly here the FFGFT side
contributes something \emph{additive}: topological winding is a homotopy
invariant; it lives in $\pi_1$ of the complex, not in the shape of its faces.
We present a self-contained, reproducible test script (companion file,
\texttt{numpy}/\texttt{scipy} only; \texttt{gudhi} optional) that
operationalises this statement: over a kNN filtration it computes the GF(2)
homology $b_1(\varepsilon)$, compares the \emph{whole} normalised Betti-curve
shape against control clouds by an energy-distance permutation test, runs a
deformation null test (persistence vs.\ drift, with optional persistent-cycle
matching), and scales the effect size with $N$. \textbf{Finding (honest,
conservative):} the \emph{scalar} total-persistence measure does not cleanly
separate HLV and is $k$-sensitive; the \emph{curve shape}, by contrast,
separates HLV significantly from all controls ($p<0.05$), sharply from
crystalline lattices (FCC, cubic, BCC) and only weakly from the generic
$\sqrt{2}$ quasicrystal. N-scaling shows that separation from crystalline
grows/stabilises (structural) while separation from random shrinks. The
discriminating feature is therefore \emph{``quasiperiodic cut-and-project
structure''}, \emph{not} the golden ratio specifically. A separate, less
circular diagnostic (longest persistent $H_1$ lifetime) additionally shows that
cut-project clouds (HLV \emph{and} the $\sqrt{2}$ quasicrystal) host a markedly
longer-lived $1$-cycle than lattices and random -- recorded as an observation,
explicitly \emph{not} oversold as ``topological protection''.

% ============================================================================
\section{Context: HLV, FFGFT, and the division of labour}
% ============================================================================

Krüger's HLV (Helix--Light--Vortex) programme~\cite{krueger276main} tests a
finite geometric substrate as a \emph{candidate object}, not as an ontology.
Its methodological strength is the pre-registered falsification
governance~\cite{krueger276gov}: positive, negative, partial and
underdetermined outcomes are kept apart on a claim-state ledger
(\emph{candidate / restricted / underdetermined / rejected / validated
candidate}); a positive result may not erase a negative residual from another
channel.

The D1--D3 audit rejected the strong claim ``HLV produces exact dodecahedral
Voronoi cells'' -- the exact 12-face channel is dominated by FCC/rhombic and
HCP-like controls. What survives is the weaker, better-supported statement of
a \emph{fresh-seed validated candidate}: a non-crystalline, mixed-face,
$\varphi$-skewed finite Voronoi cell ensemble. The resulting reformulation of
the vortex axiom reads, in essence: vortex and boundary modes are supported by
the edge, face-boundary, loop, and defect-interface structure of a mixed-face
complex -- not by an exact cell shape.

\begin{observation*}
Topological winding requires no Platonic cell. It requires closed graph
cycles, phase continuity, and holonomy-carrying loops. These structures live
in $\pi_1$ of the complex and are invariant under the cell shape.
\end{observation*}

This is also the FFGFT position: on the compact $T^4$ the relevant invariants
are holonomy classes of $\pi_1(T^4)=\mathbb{Z}^4$, deliberately metric- and
cell-independent (cf.\ Doc.~270~\cite{dok270}, Doc.~271~\cite{dok271}). The
division of labour remains complementary, not overlapping: HLV as the
boundary-pattern residual test, FFGFT as the scale-law / effective-$L_\xi$
test (Doc.~275~\cite{dok275}).

% ============================================================================
\section{The test method}
% ============================================================================

\subsection{Complex and boundary operator}
Instead of counting Voronoi faces, we build the clique complex of the
symmetric kNN graph (the same graph object as in Krüger's runs, $k$
adjustable) and read the cycle/holonomy sector off the GF(2) homology
\[
  H_1 = \ker \partial_1 / \operatorname{im} \partial_2,
  \qquad
  b_1 = \dim\ker\partial_1 - \operatorname{rank}\partial_2
      = (E - \operatorname{rank}\partial_1) - \operatorname{rank}\partial_2 .
\]
Here $\operatorname{rank}\partial_1 = N - (\#\text{components})$ (union-find),
and $\operatorname{rank}\partial_2$ is computed over GF(2). $b_1$ is the
cell-shape-independent winding/holonomy capacity of the complex.

\paragraph{Language precision.} $H_1$ is shape-independent \emph{within a fixed
complex}; it still depends on $k$ and the metric. Hence (i) the $k$-sweep and
(ii) the comparison of the \emph{whole} curve shape rather than a single
``tessellation-free'' number.

\subsection{Filtration and curve shape}
A dense kNN complex is nearly contractible; the holonomy structure becomes
visible only through a radius filtration. We evaluate $b_1(\varepsilon)$ on a
grid $\varepsilon/d_0 \in \{1.1,\,1.35,\,1.6,\,1.85,\,2.1,\,2.35\}$ ($d_0$ =
median nearest-neighbour distance) and compare the \emph{normalised} curve
$\hat b_1 = b_1/\sum b_1$ -- the shape, not the area. The scalar total
persistence $\int b_1\,\mathrm{d}\varepsilon$ is kept only as a coarse
single-number summary.

\subsection{Statistics and controls}
Each cloud is promoted to an ensemble over seeds (the same
``HLV-as-an-ensemble'' discipline as Krüger's Run~E). Scalars are compared by a
permutation test, curve shapes by an energy-distance permutation test.
Controls: FCC, simple cubic, BCC (crystalline); jittered cubic, uniform
random, Poisson-disc (disordered, against density artefacts); generic
$\sqrt{2}$ cut-and-project (quasicrystal without the golden ratio).

\subsection{Deformation null test}
We morph HLV $\to$ control continuously, $x(t)=(1-t)\,x_{\mathrm{HLV}}
+ t\,x_{\mathrm{control}}$, and track total persistence. With \texttt{gudhi}
present we additionally track the \emph{longest} $H_1$ lifetime (Rips
persistence) -- far less circular, since a single long bar is not forced to
$0$. Caveat: a target with $b_1\approx 0$ (e.g.\ FCC) forces the endpoint to
$0$; clean cycle matching is only possible via persistent homology.

\subsection{N-scaling and pre-registered consistency check}
The effect size (L1 distance of mean curve shapes, HLV vs.\ the pooled
crystalline profile) is scaled with $N$: stabilising $\Rightarrow$ structural,
vanishing $\Rightarrow$ finite-size artefact. The \texttt{-{}-falsify} switch
implements a \emph{pre-registered consistency check} (fixed $N,k,p$;
crystalline controls only) -- explicitly \emph{not} a Popperian falsification
test: a negative result weakens the claim but does not logically falsify it.

% ============================================================================
\section{Results}
% ============================================================================

All numbers are from a real run of the companion script ($N=240$, $k=14$,
seed base $20260613$, $5$ seeds).

\paragraph{Scalar (Check 1/1b).} The total-persistence measure is unstable and
$k$-sensitive: the simple-cubic control dominates ($\sim\!130$ vs.\ HLV
$\sim\!50$), and over $k\in\{10,14,18\}$ HLV lies above the control band in
$0/3$ cases. Claim-state: \emph{underdetermined}. This is precisely why the
scalar is discarded and the shape compared.

\paragraph{Curve shape (Check 1c).} The energy-distance test separates HLV
from \emph{all} controls ($p<0.05$):

\begin{center}
\begin{tabular}{lrr}
\hline
Control & energy distance & $p$ \\
\hline
jittered cubic        & $2.81$ & $0.0015$ \\
uniform random        & $2.79$ & $0.0060$ \\
Poisson-disc          & $2.71$ & $0.0030$ \\
FCC (rhombic-dodeca.) & $1.93$ & $0.0075$ \\
cubic                 & $1.67$ & $0.0040$ \\
BCC                   & $1.11$ & $0.0060$ \\
generic QC ($\sqrt2$) & $0.15$ & $0.0085$ \\
\hline
\end{tabular}
\end{center}

\noindent The separation is sharp against crystalline lattices and very weak
(smallest distance $0.15$) against the generic $\sqrt{2}$ quasicrystal.

\paragraph{Longest $H_1$ lifetime (with \texttt{gudhi}, max edge $=2\,d_0$).}

\begin{center}
\begin{tabular}{lrr}
\hline
Cloud & raw & normalised ($/d_0$) \\
\hline
HLV (golden cut-project)    & $2.07$ & $0.87$ \\
generic QC ($\sqrt2$)       & $2.01$ & $0.85$ \\
cubic                       & $0.41$ & $0.41$ \\
uniform random              & $0.25$ & $0.43$ \\
BCC                         & $0.13$ & $0.15$ \\
FCC                         & $0.00$ & $0.00$ \\
\hline
\end{tabular}
\end{center}

\paragraph{N-scaling.} The mean-shape L1 distance HLV-vs-crystalline
grows/stabilises with $N$ ($\approx 0.56 \to 0.67 \to 1.61 \to 1.51$ for
$N=200/400/600/800$) -- structural, not a pure finite-size artefact;
HLV-vs-random shrinks accordingly.

\paragraph{Deformation and consistency check.} In the deformation null test
both total persistence and longest $H_1$ lifetime fall; at this $N$ there is
\emph{no} robust plateau (status: underdetermined -- honest). The
pre-registered consistency check (fixed $N,k,p<0.01$; FCC/cubic/BCC only)
returns \emph{PASS}: HLV separates from every crystalline control.

\begin{figure}[htbp]
  \centering
  \includegraphics[width=\linewidth]{276_homology_figure.png}
  \caption{Four-panel overview from the companion script (v3). A: normalised
  Betti curves $b_1(\varepsilon)$ -- shape, not magnitude. B: curve-shape
  energy distance HLV vs.\ controls. C: deformation null test (total
  persistence and, with \texttt{gudhi}, longest $H_1$ lifetime). D: N-scaling
  of the effect size (stable $=$ structural). Finite diagnostic; no claim about
  physical spacetime microstructure.}
  \label{fig:276homology}
\end{figure}

% ============================================================================
\section{Discussion and claim-state}
% ============================================================================

The defensible statement is precise and conservative:
\begin{enumerate}
  \item HLV separates \emph{structurally} from crystalline lattices
        (FCC/cubic/BCC) via the $b_1(\varepsilon)$ curve shape; the effect
        grows with $N$ (claim-state: \emph{candidate}, shape-based
        $k$-robust).
  \item HLV is nearly indistinguishable from the generic $\sqrt{2}$
        cut-and-project quasicrystal on \emph{every} metric (shape energy
        distance $0.15$; longest cycle $2.07$ vs.\ $2.01$). The discriminating
        feature is therefore \emph{``quasiperiodic cut-and-project
        structure''}, \emph{not} the golden ratio specifically.
  \item cut-project structures host a markedly longer-lived $H_1$ cycle
        ($\sim\!2\,d_0$) than lattices and random ($\le 0.4\,d_0$). This is a
        real observation of the less circular persistence diagnostic; it is
        \emph{not} interpreted as ``topological protection'' -- that would be
        the same overclaim with the opposite sign.
\end{enumerate}

Methodologically the script confirms the claim-state discipline of the HLV
programme: the obvious scalar measure fails, and progress lies in the
shape-based, ensemble- and permutation-supported comparison. Stronger
statements require larger $N$, more seeds, exact repository adjacency, and
independent reimplementation.

\paragraph{Reproducibility.} The companion script
(\texttt{Dok276\_Skripte/ffgft\_hlv\_homology\_check\_v3.py}) depends only on
\texttt{numpy}/\texttt{scipy}; \texttt{gudhi} and \texttt{matplotlib} are
optional (true persistence diagrams and the figure, respectively).
Deliberately no hard dependency -- anyone on the IPI list can reproduce the
result without a build step.

\paragraph{Acknowledgement.} Marcel Krüger (HLV G-Lattice, falsification
governance) for the open exchange and the pre-registered test architecture on
which this document builds.

% ============================================================================
\begin{thebibliography}{9}

\bibitem{krueger276main}
M.~Krüger,
\emph{Operational Falsification Pathways for the HLV G-Lattice:
Weighted Spectral Forks, Fresh-Seed Cell-Shape Validation, Native 6D
Operators, and Process-Tensor Bridge Tests},
Zenodo (2026), \url{https://doi.org/10.5281/zenodo.20668125}.

\bibitem{krueger276gov}
M.~Krüger,
\emph{Operational Falsification Governance: Residual Preservation,
Claim-State Ledgers, and Revision-Admissible Scientific Frameworks},
Zenodo (2026), \url{https://doi.org/10.5281/zenodo.20629842}.

\bibitem{dok275}
J.~Pascher, \emph{Doc.\,275: From $\xi$ to $\varphi$ -- Recursive scale
hierarchy and spectral gap structure}, FFGFT corpus (2026).

\bibitem{dok271}
J.~Pascher, \emph{Doc.\,271: FFGFT compared with HLV}, FFGFT corpus (2026).

\bibitem{dok270}
J.~Pascher, \emph{Doc.\,270: Dimensional change in FFGFT
($T^4 \to T^0$)}, FFGFT corpus (2026).

\bibitem{edelsbrunner}
H.~Edelsbrunner and J.~Harer, \emph{Computational Topology: An
Introduction}, American Mathematical Society (2010).

\bibitem{gretton}
A.~Gretton et al., \emph{A kernel two-sample test},
Journal of Machine Learning Research \textbf{13}, 723--773 (2012).

\end{thebibliography}
